%%%%%%%%%%%%%%%%%%%% CITIS-Final-Paper-Adaptado.tex %%%%%%%%%%%%%%%%%%%%%%
%
%  MARCADORES \todo{...}: senalan datos que deben recalcularse a partir
%  de las corridas reales, donde depende de las corridas por semilla (rangos de Friedman,
%  p-valores, Nemenyi) se dejo el marcador para que ustedes lo
%  recalculen con el pipeline de estadistica.
% =====================================================================

\documentclass[graybox]{SNmult}

\usepackage{type1cm}
\usepackage{makeidx}
\usepackage{graphicx}
\usepackage{multicol}
\usepackage[bottom]{footmisc}
\usepackage{newtxtext}
\usepackage[varvw]{newtxmath}
\usepackage{amsmath}
\let\Bbbk\undefined % Resolves amssymb \Bbbk collision
\usepackage{amssymb}
\usepackage{booktabs}
\usepackage{algorithm}
\usepackage{algpseudocode}
\usepackage{url}
\usepackage{xcolor}

\makeindex             % como en la plantilla: usar el estilo svind.ist
                       % con su programa makeindex

% Marcador visible para pendientes de recalculo/edicion antes de enviar.
% Buscar "\todo{" en el fuente para ubicar todos los pendientes.
\newcommand{\todo}[1]{\textcolor{red}{\textbf{[PENDIENTE: #1]}}}

\begin{document}
% \motto{...}   % opcional: la plantilla admite un motto por capitulo;
                % se deja sin usar salvo que los autores quieran uno.
\title*{Capacitated School Bus Routing via Clustering and
Network-Based Optimization: A Comparative Benchmark of Six Strategies
on a Real Road Network}
\titlerunning{Capacitated School Bus Routing: A Comparative Benchmark}

% --- Bloque de autoria camera-ready ---------------------------------
\author{Roy Sandoval\orcidID{0009-0009-0824-2390} \and
Sara R\'ua\orcidID{0009-0002-7170-7479} \and
Mariana Arbel\'aez\orcidID{0009-0005-3642-2760} \and
Naomi Chow\orcidID{0009-0003-0356-3167} \and
Daniel L\'opez\orcidID{0009-0001-9911-2904} \and
Roberto Hincapi\'e\orcidID{0000-0003-2525-1247}}
\authorrunning{Sandoval et al.}
\institute{Roy Sandoval \at Ingenier\'ia en Ciencia de Datos,
Universidad Pontificia Bolivariana, Medell\'in, Colombia
\email{roy.sandoval@upb.edu.co}
\and Sara R\'ua \at Ingenier\'ia en Ciencia de Datos,
Universidad Pontificia Bolivariana, Medell\'in, Colombia
\email{sara.ruag@upb.edu.co}
\and Mariana Arbel\'aez \at Ingenier\'ia en Ciencia de Datos,
Universidad Pontificia Bolivariana, Medell\'in, Colombia
\email{mariana.arbelaezs@upb.edu.co}
\and Naomi Chow \at Ingenier\'ia en Ciencia de Datos,
Universidad Pontificia Bolivariana, Medell\'in, Colombia
\email{naomi.chow@upb.edu.co}
\and Daniel L\'opez \at Ingenier\'ia en Ciencia de Datos,
Universidad Pontificia Bolivariana, Medell\'in, Colombia
\email{daniel.lopezlo@upb.edu.co}
\and Roberto Hincapi\'e \at Ingenier\'ia en Ciencia de Datos,
Universidad Pontificia Bolivariana, Medell\'in, Colombia
\email{roberto.hincapie@upb.edu.co}}

\maketitle

\abstract*{The School Bus Routing Problem (SBRP) is an NP-hard
combinatorial optimization problem that generalizes the capacitated
vehicle routing problem and underpins a service used daily by millions
of students. Although the literature offers many heuristic and
metaheuristic strategies, head-to-head comparisons under identical,
realistic conditions remain scarce, which makes it difficult for
practitioners to choose a method. We present a unified benchmark that
evaluates six clustering-, covering- and search-based routing
strategies: capacitated k-medoids, a multi-objective sectorial
clustering, a zonal capacitated vehicle routing problem with time
windows (CVRPTW), two set-covering integer programs, and an adaptive
large-neighbourhood search (ALNS), on identical scenarios generated
from the real OpenStreetMap road network of Medell\'in (Colombia). All
methods share a common solver interface, directed shortest-path
distances, a single depot, and a bus capacity of twenty, and are
measured on fleet distance, computation time, fleet size, coverage and
cluster cohesion across student densities $N\in\{50,100,200,400\}$ with
eleven seeds each (264 runs). Capacitated k-medoids and sectorial
clustering are statistically tied for the shortest fleet distance
(Wilcoxon $p=0.36$), while the set-covering formulations trade route
quality for speed, at different points of that trade-off because they
optimize different objectives. Under an identical wall-clock budget the
ALNS attains the shortest fleet distance at every density, winning all
$44$ scenarios against both clustering pipelines and taking first place on
the distance ranking (mean rank $1.02$ of $6$), while calibration against
published CVRP instances puts it within $0.22\%$ of the known optima on
average. It reaches those routes with markedly \emph{less} cohesive
clusters than the clustering pipelines, which indicates that spatial
compactness is a by-product of good routes rather than their cause. We
distill the trade-offs into practical guidance for SBRP solver selection.
\keywords{School Bus Routing Problem $\cdot$ Capacitated clustering $\cdot$
Vehicle routing $\cdot$ Set covering $\cdot$ Road-network optimization.}}

\abstract{The School Bus Routing Problem (SBRP) is an NP-hard
combinatorial optimization problem that generalizes the capacitated
vehicle routing problem and underpins a service used daily by millions
of students. Although the literature offers many heuristic and
metaheuristic strategies, head-to-head comparisons under identical,
realistic conditions remain scarce, which makes it difficult for
practitioners to choose a method. We present a unified benchmark that
evaluates six clustering-, covering- and search-based routing
strategies: capacitated k-medoids, a multi-objective sectorial
clustering, a zonal capacitated vehicle routing problem with time
windows (CVRPTW), two set-covering integer programs, and an adaptive
large-neighbourhood search (ALNS), on identical scenarios generated
from the real OpenStreetMap road network of Medell\'in (Colombia). All
methods share a common solver interface, directed shortest-path
distances, a single depot, and a bus capacity of twenty, and are
measured on fleet distance, computation time, fleet size, coverage and
cluster cohesion across student densities $N\in\{50,100,200,400\}$ with
eleven seeds each (264 runs). Capacitated k-medoids and sectorial
clustering are statistically tied for the shortest fleet distance
(Wilcoxon $p=0.36$), while the set-covering formulations trade route
quality for speed, at different points of that trade-off because they
optimize different objectives. Under an identical wall-clock budget the
ALNS attains the shortest fleet distance at every density, winning all
$44$ scenarios against both clustering pipelines and taking first place on
the distance ranking (mean rank $1.02$ of $6$), while calibration against
published CVRP instances puts it within $0.22\%$ of the known optima on
average. It reaches those routes with markedly \emph{less} cohesive
clusters than the clustering pipelines, which indicates that spatial
compactness is a by-product of good routes rather than their cause. We
distill the trade-offs into practical guidance for SBRP solver selection.
\keywords{School Bus Routing Problem $\cdot$ Capacitated clustering $\cdot$
Vehicle routing $\cdot$ Set covering $\cdot$ Road-network optimization.}}

% =====================================================================
\section{Introduction}
% =====================================================================
The School Bus Routing Problem\index{School Bus Routing Problem} (SBRP) consists of planning a fleet of
buses that collects students from a set of stops and delivers them to
their school while respecting operational constraints such as vehicle
capacity, maximum riding time and bell-time windows~\cite{Park2010TheReview}.
It is a generalization of the capacitated vehicle routing problem
(CVRP), which itself descends from the truck-dispatching problem first
formalized by Dantzig and Ramser~\cite{DantzigRamser1959Truck}, and it
is therefore NP-hard. The intractability is also combinatorial: a bus
visiting $n_c$ stops admits $n_c!$ visiting orders, and the partition
of students across buses must be chosen jointly with those orders, so
naive enumeration grows factorially in the number of stops. Exact
methods are therefore infeasible at realistic $n$, and practical SBRP
solvers are overwhelmingly heuristic. The survey by Park and
Kim~\cite{Park2010TheReview}, together with the more recent review of
Ellegood et al.~\cite{Ellegood2020SchoolDirections}, shows that the
field is fragmented into many problem variants and solution paradigms
whose relative merits are rarely measured on common ground.

A classical and still dominant family of heuristics follows the
\emph{cluster-first, route-second}\index{cluster-first, route-second}
paradigm: students are first
grouped into capacity-feasible clusters, and a Traveling Salesman
Problem (TSP) is then solved within each cluster, for which highly
effective heuristics exist, notably Croes' 2-opt local
search~\cite{Croes1958TwoOpt} and Helsgaun's implementation of the
Lin--Kernighan heuristic~\cite{HelsgaunAnHeuristic}. The clustering
step can be guided by geometry~\cite{CakirRevisitingProblem}, by
classical partitioning algorithms~\cite{Comert2018AStudy}, or by
multiple and often conflicting objectives such as simultaneously
minimizing the number of buses and the students' riding
time~\cite{Corberan2002HeuristicObjectives}. An orthogonal family
formulates route selection directly as a set-covering\index{set covering}
integer program.
A third relies on population-based metaheuristics such as genetic,
annealing or nature-inspired algorithms~\cite{Baranwal2016DAVRPTW,Hatamlou2018BlackHole};
Section~\ref{sec:alns} describes the adaptive large-neighbourhood
search from this family that we include as the sixth strategy in the
present benchmark.

The difficulty for a practitioner is that these strategies are
typically published on different instances, with different distance
models (Euclidean vs.\ road network), different objectives and
different computational budgets, so their reported numbers cannot be
compared directly. This paper addresses that gap. Our contributions
are:

\begin{itemize}
\item A \textbf{unified benchmark harness} exposing six SBRP
strategies through one interface and evaluating them on \emph{identical}
scenarios built from the real road network of Medell\'in and the Valle
de Aburr\'a, under the same capacity, distance model and metrics.
\item A \textbf{rigorous description} of each strategy: a formal
decision-variable model, and worst-case complexity checked against
\emph{measured} per-stage cost. We also make explicit where two
strategies that look comparable in fact optimize different objectives,
so that a shared metric (e.g.\ distance) does not silently conflate
objective choice with implementation quality.
\item A \textbf{statistically grounded comparison} across four
densities and eleven seeds (264 runs), using Friedman and Wilcoxon
tests with a Nemenyi post-hoc to separate genuine differences from
noise.
\item \textbf{Practical guidance}: capacitated k-medoids and sectorial
clustering are statistically tied for the best fleet distance, while
set-covering programs occupy the speed end of the speed/quality
trade-off. Section~\ref{sec:discussion} explains why, and
Section~\ref{sec:alns} describes the adaptive large-neighbourhood
search (ALNS) we add as a sixth, locality-aware baseline so that a
population-based metaheuristic is measured on the same scenarios,
metrics and solver interface as the other five strategies.
\end{itemize}

\paragraph{Relevance to data-driven urban planning.}
This work addresses \emph{smart mobility and data-driven urban
planning}: it decides school transport from the real street topology of
a metropolitan area, and re-targets to another city by swapping the
graph and service polygons. Total fleet distance is exactly the
vehicle-kilometres travelled (VKT) that drives fuel use, emissions and
congestion, so the spread between solvers is not merely algorithmic. At
$N=400$ the worst strategy (Set Cover, per-child, $670$~km) commits about
$1.56\times$ the VKT of the best (ALNS, $428$~km) for the same students,
buses and network; the $32$~km that ALNS saves over the best clustering
pipeline alone is $7\%$ of the daily fleet distance, recovered purely by
changing solver. We stop at VKT
deliberately: converting it to emissions needs fleet-specific factors
we lack for the study area, so that step, and validation against
municipal operations, are future work rather than an estimate we
invent~\cite{Sabet2022GreenVRP}.

% =====================================================================
\section{Related Work}\label{sec:related}
% =====================================================================
\subsection{The School Bus Routing Problem}
The SBRP has been studied for almost five decades. Park and
Kim~\cite{Park2010TheReview} decompose it into five sub-problems (data
preparation, bus-stop selection, route generation, route scheduling
and bell-time adjustment) and observe that most contributions tackle
one or two of these in isolation. The updated survey of Ellegood et
al.~\cite{Ellegood2020SchoolDirections} reviews 64 newer publications
and confirms that single-school, capacity-constrained route generation
remains the most studied core, and that heuristic and metaheuristic
methods dominate exact ones in practice. The most recent synthesis, by
Wang et al.~\cite{Wang2025SBRPReview}, covers the literature through
2024 and records a shift from single-objective cost minimization
towards multi-objective formulations that add heterogeneous fleets,
service equity and dynamic constraints, alongside marked growth in
metaheuristic and learning-based methods; it also names what motivates
this paper, namely that the proliferation of variants has outpaced
comparative evaluation. Corber\'an et
al.~\cite{Corberan2002HeuristicObjectives} emphasize the inherently
\emph{multi-objective} nature of the problem: minimizing the number of
buses and minimizing student travel time are conflicting goals, and a
good solver must navigate that trade-off rather than optimize a single
scalar. Our benchmark adopts the single-school, capacitated,
route-generation variant and reports several of these competing
objectives side by side.

\subsection{Cluster-first, route-second}
Clustering was central to early VRP algorithms and has seen renewed
interest. Cakir et al.~\cite{CakirRevisitingProblem} revisit
cluster-first, route-second with a shape-based clustering that forms
clusters of a desired geometry relative to the depot, which is the
conceptual ancestor of our sectorial strategy. Comert et
al.~\cite{Comert2018AStudy} solve a capacitated VRP by first
partitioning customers with k-means, k-medoids or random clustering
and then routing each cluster, and report that medoid-based clustering
yields the most robust routes, a finding our results corroborate on a
real road network. The k-medoids/PAM family itself originates with
Kaufman and Rousseeuw~\cite{KaufmanRousseeuw1990PAM}, and we use the
FasterPAM variant for initialization. Geometry- and curve-based
clustering continues to advance, e.g.\ the sub-trajectory set-cover
formulation of Akitaya et al.~\cite{Akitaya2021Subtrajectory}, which
also links clustering to set covering, the two paradigms we contrast.

\subsection{Routing and metaheuristics}
Within a cluster, route construction reduces to the TSP. The savings
heuristic of Clarke and Wright~\cite{ClarkeWright1964Savings} and local
search such as 2-opt~\cite{Croes1958TwoOpt} are the classical building
blocks, while the Lin--Kernighan heuristic and its
implementations~\cite{HelsgaunAnHeuristic} remain the state of the art
for high-quality tours. Modern solvers wrap these in metaheuristics;
Google OR-Tools~\cite{PerronFurnon2023ORTools}, which we use for the
routing stage, drives local search with Guided Local
Search~\cite{Voudouris1999GLS}. Population- and physics-inspired
metaheuristics are also widely applied to routing: deterministic
annealing for the VRP with time windows~\cite{Baranwal2016DAVRPTW},
the black-hole algorithm for the TSP~\cite{Hatamlou2018BlackHole}, and
genetic algorithms~\cite{Holland1992Adaptation}. More recently,
learning-based approaches such as graph convolutional networks have
been proposed for the TSP~\cite{Joshi2019GCN,Dahiya2018TSPReview}, and
neural combinatorial optimization has matured into a distinct family of
VRP solvers~\cite{Wu2024NCOSurvey}. A parallel line makes
sustainability an explicit objective, formulating emissions- and
energy-aware ``green'' VRP variants~\cite{Sabet2022GreenVRP}. The
present benchmark compares six clustering-, covering- and search-based
strategies; adaptive large-neighbourhood search (ALNS) is included as
the population-based, locality-aware representative of the
metaheuristic family, described in Section~\ref{sec:alns}, though a
single ALNS configuration does not claim to bound that broader family.

% =====================================================================
\section{Methods}\label{sec:methods}
% =====================================================================
This section presents the methods in three parts: the formal problem
statement (Section~\ref{sec:problem}), the six solution strategies
with their models and complexity, including the adaptive
large-neighbourhood search (ALNS)
(Sections~\ref{sec:framework}--\ref{sec:alns}), and the experimental
protocol (Section~\ref{sec:setup}).

\subsection{Problem formulation}\label{sec:problem}
\subsubsection{Notation and distances.}
Let the road network be a directed weighted graph $G=(V,E)$ with arc
weights $w_e>0$. A single school (depot) $o\in V$ and a set of
$n$ student stops $\mathcal{C}=\{1,\dots,n\}\subset V$ are given,
together with a homogeneous fleet of buses of capacity $C$. For any
pair $i,j\in V$ the road distance is the shortest-path cost
\begin{equation}
d_{ij}=\min_{p\in\mathcal{P}(i,j)}\sum_{e\in p} w_e ,
\end{equation}
where $\mathcal{P}(i,j)$ is the set of directed paths from $i$ to $j$;
distances are in general asymmetric ($d_{ij}\neq d_{ji}$). Clustering
algorithms that assume a metric use the symmetrized matrix
\begin{equation}
D^{\mathrm{sym}}_{ij}=\tfrac{1}{2}\,(d_{ij}+d_{ji}) ,
\end{equation}
whereas the routing stage uses the directed distances $d_{ij}$.

\subsubsection{Objective and constraints.}
The SBRP we address is a single-depot CVRP. A solution is a set of
closed routes, one per bus, each starting and ending at $o$. The
number of buses is lower-bounded by
$k=\lceil n/C\rceil$. The objective is to minimize the total fleet
distance
\begin{equation}
L_{\mathrm{total}}=\sum_{c}\sum_{(i,j)\in r_c} d_{ij}\,,
\label{eq:obj}
\end{equation}
subject to: (i)~every student is served by exactly one bus (full
coverage); (ii)~no route carries more than $C$ students; and
(iii)~every route starts and ends at the depot. Because solving this
CVRP exactly is intractable, we adopt the cluster-first, route-second
decomposition into a capacitated assignment sub-problem followed by
$k$ independent TSPs.

Constraint~(i) is hard for five of the six strategies, including ALNS
(Section~\ref{sec:alns}), which repairs every destroyed solution back
to full coverage before it is accepted. The zonal
CVRPTW is the exception: it also enforces a route-duration limit, which
can conflict with coverage on peripheral stops, so it relaxes~(i) to
$\sum_{c} x_{ic}+u_i=1$ with $u_i\in\{0,1\}$ marking student $i$
unserved and adds $\Pi\sum_i u_i$ to the objective, for $\Pi$ large
enough that a drop is worse than any feasible detour
(Section~\ref{sec:cvrptw}). Coverage is reported as a first-class
metric so this relaxation stays visible in every comparison.

\subsubsection{Capacitated assignment sub-problem.}
Given current cluster representatives (medoids)
$M=\{m_1,\dots,m_k\}$, the assignment of students to clusters is the
binary program
\begin{align}
\min \quad & \sum_{i=1}^{n}\sum_{c=1}^{k} D^{\mathrm{sym}}_{i,m_c}\,x_{ic}
\label{eq:assign-obj}\\
\text{s.t.}\quad & \sum_{c=1}^{k} x_{ic}=1 \quad \forall i, \label{eq:assign-cover}\\
& \sum_{i=1}^{n} x_{ic}\le C \quad \forall c, \label{eq:assign-cap}\\
& x_{ic}\in\{0,1\},
\end{align}
where $x_{ic}=1$ iff student $i$ joins cluster $c$. As the constraint
matrix of (\ref{eq:assign-obj})--(\ref{eq:assign-cap}) is totally
unimodular, the program is solved to optimality as a minimum-cost flow
in a four-layer network (source $\to$ students $\to$ medoids $\to$ sink),
avoiding a general integer solver. Each cluster's medoid is then
recomputed as the within-cluster $1$-median
$m_c=\arg\min_{h\in S_c}\sum_{i\in S_c} D^{\mathrm{sym}}_{ih}$, and the
two steps alternate until the medoids stabilize.

\subsubsection{Per-cluster routing sub-problem.}
For cluster $c$ with node set $V_c=S_c\cup\{o\}$ and binary arc
variables $y_{ij}$, the route is the TSP tour
\begin{align}
\min \quad & \sum_{(i,j)\in V_c\times V_c} d_{ij}\,y_{ij}\\
\text{s.t.}\quad & \textstyle\sum_{j\neq i} y_{ij}=1,\;\;
\sum_{i\neq j} y_{ij}=1 \quad \forall i,j\in V_c,\\
& \textstyle\sum_{i\in S}\sum_{j\in S} y_{ij}\le |S|-1
  \quad \forall\, S\subsetneq V_c,\; 2\le|S|\le|V_c|-1,
  \label{eq:dfj}\\
& y_{ij}\in\{0,1\},
\end{align}
where (\ref{eq:dfj}) is the Dantzig--Fulkerson--Johnson
subtour-elimination family: forbidding any proper subset $S$ from
carrying $|S|$ internal arcs rules out cycles that do not span $V_c$.
Having exponentially many members, it is not enumerated: OR-Tools
maintains a single connected circuit by construction and improves it
with Guided Local Search~\cite{Voudouris1999GLS} under a fixed budget,
so (\ref{eq:dfj}) defines the optimum we approximate, not the algorithm.

\subsubsection{Set-covering formulation.}\label{sec:setcover-formulation}
An alternative that bypasses explicit clustering enumerates a pool
$\mathcal{R}$ of capacity-feasible candidate routes and selects a
subset. Minimizing the number of buses gives the set-cover program
\begin{equation}
\min \sum_{j\in\mathcal{R}} x_j \quad\text{s.t.}\quad
\sum_{j:\,i\in R_j} x_j \ge 1\;\;\forall i\in\mathcal{C},\quad
x_j\in\{0,1\},
\label{eq:setcover}
\end{equation}
while replacing the objective by $\sum_j \mathrm{cost}_j\,x_j$ (the
route length) minimizes distance instead. Both variants are solved
with a branch-and-cut integer solver. \textbf{This objective choice is
not incidental}: the two set-covering strategies we implement
(Section~\ref{sec:setcover-two}) instantiate different sides of
Eq.~(\ref{eq:setcover}), one minimizes fleet size and the other
minimizes distance, so any later comparison of the two on a single
metric such as distance mixes two effects that our design does not
separate: how the candidate pool $\mathcal{R}$ is built, and which
objective is being optimized. We return to this point in
Section~\ref{sec:setcover-two} and again in Section~\ref{sec:discussion}.

\subsection{Common benchmark framework}\label{sec:framework}
All strategies are implemented as adapters behind a single
\texttt{Solver} interface: each receives a \texttt{Scenario} (the same
road graph, student locations, depot and capacity) and returns a
\texttt{Solution} (a list of closed routes and, optionally, cluster
labels). This guarantees that the strategies are evaluated on
byte-identical inputs and scored by the same metric code, removing
instance and implementation bias from the comparison. The framework is
written in Python and relies on OSMnx~\cite{Boeing2017OSMnx} for
road-network extraction, NetworkX for graph algorithms and minimum-cost
flow, OR-Tools~\cite{PerronFurnon2023ORTools} for TSP/CVRPTW search,
FasterPAM~\cite{KaufmanRousseeuw1990PAM} for medoid initialization, a
branch-and-cut integer solver (CBC) for the set-covering programs, and
scikit-learn~\cite{Pedregosa2011Sklearn} for the silhouette score. The
scenario generator is deterministic given the pair $(N,\text{seed})$, so
every instance is reproducible. This same interface is what lets the
ALNS adapter (Section~\ref{sec:alns}) enter the comparison without
touching the scenarios, metrics or any other strategy.

\subsection{Capacitated k-medoids + OR-Tools (baseline)}
The baseline instantiates the cluster-first, route-second
decomposition exactly as formalized above (Algorithm~\ref{alg:kmedoids}):
FasterPAM provides initial medoids, the capacitated assignment is
solved by minimum-cost flow, medoids are recomputed, and the loop
iterates to convergence before each cluster is routed with OR-Tools.

\begin{algorithm}[t]
\caption{Capacitated k-medoids + per-cluster TSP}
\label{alg:kmedoids}
\footnotesize
\begin{algorithmic}[1]
\Require symmetrized distances $D^{\mathrm{sym}}$, directed distances $d$,
capacity $C$, depot $o$, iteration cap $T$
\State $k \gets \lceil n/C\rceil$
\State $M \gets \textsc{FasterPAM}(D^{\mathrm{sym}},k)$ \Comment{initial medoids}
\Repeat
  \State $x \gets \textsc{MinCostFlowAssign}(D^{\mathrm{sym}},M,C)$
         \Comment{Eqs.~(\ref{eq:assign-obj})--(\ref{eq:assign-cap})}
  \ForAll{clusters $S_c$}
     \State $m_c \gets \arg\min_{h\in S_c}\sum_{i\in S_c} D^{\mathrm{sym}}_{ih}$
  \EndFor
\Until{$M$ unchanged \textbf{or} $T$ iterations}
\ForAll{clusters $S_c$}
  \State $r_c \gets \textsc{OrToolsTSP}(S_c\cup\{o\},d)$ \Comment{Guided Local Search}
\EndFor
\State \Return $\{r_c\}$
\end{algorithmic}
\end{algorithm}

\subsection{Sectorial multi-objective clustering}
The sectorial strategy replaces the purely distance-driven assignment
by a geometric one inspired by shape-based
clustering~\cite{CakirRevisitingProblem}. Each student is scored by a
convex combination of its angular position and its radial distance
with respect to the school,
\begin{equation}
s_i=\alpha\,\frac{\theta_i+\pi}{2\pi}+\beta\,\hat d_i,
\qquad \alpha+\beta=1,
\label{eq:sectorial-score}
\end{equation}
where $\theta_i=\operatorname{atan2}(y_i-y_o,\,x_i-x_o)$ and
$\hat d_i$ is the min--max-normalized distance to the depot
($\alpha=0.6,\beta=0.4$ by default). An angular sweep packs students
into capacity-bounded sectors, which are then refined with the same
minimum-cost-flow reassignment as the baseline, optionally penalizing
moves that break the initial sweep. The result is a set of compact,
wedge-shaped clusters that radiate from the school.

\subsection{Zonal CVRPTW}\label{sec:cvrptw}
This strategy partitions the metropolitan area into six fixed
geographic zones (communes and municipalities of the Valle de
Aburr\'a, by point-in-polygon test) and solves an independent CVRPTW
per zone with OR-Tools. Beyond capacity it enforces a time dimension:
travel time plus a $90$-second boarding time per stop, a maximum route
duration of $75$ minutes, and a span-balancing penalty that
discourages long routes. Unserved stops are allowed at a high penalty,
so the method can leave a small fraction of students uncovered when the
duration limit binds.

Formally, with $Z$ indexing zones, $\mathcal{C}_z$ the students inside
zone $z$, binary arc variables $y^z_{ij}$, drop indicators
$u_i\in\{0,1\}$ and arrival times $a_i\ge0$, each zone solves
\begin{align}
\min \;\; & \sum_{z\in Z}\Big[\sum_{(i,j)} d_{ij}\,y^z_{ij}
  + \gamma\,\sigma_z\Big] + \Pi\sum_{i\in\mathcal{C}} u_i
\label{eq:cvrptw}\\
\text{s.t.}\;\; & \textstyle\sum_{j} y^z_{ij} = 1-u_i,\quad
\sum_{i\in r} q_i \le C,\quad
a_j \ge a_i + \tau_{ij} + b,
\nonumber
\end{align}
plus $a_{\text{end}}-a_{\text{start}}\le T_{\max}$ per route, with
$\tau_{ij}$ directed travel time, $b=90$~s boarding, $T_{\max}=75$~min,
$\sigma_z$ the duration span penalized by $\gamma$, and $\Pi$ the drop
penalty of Section~\ref{sec:problem}. The constraint
$\sum_j y^z_{ij}=1-u_i$ is the soft-coverage relaxation: a student is
visited once or dropped at cost $\Pi$, and since $\Pi$ dominates any
routing saving, drops arise only where the duration limit admits no
feasible insertion. With a single bell time the time windows reduce to
that limit, which is what binds here.

\subsection{Set-covering integer programs}\label{sec:setcover-two}
Two variants instantiate the set-cover model of
Eq.~(\ref{eq:setcover}), and, as flagged above, they instantiate it
\emph{differently}: they build different candidate pools and they
optimize different objectives, so they are two distinct policies, not
two implementations of the same policy. The \emph{per-child} variant
builds one candidate route per student by greedily packing its $C-1$
nearest neighbours, then selects the minimum number of routes covering
all students, i.e.\ it minimizes fleet size. The \emph{grid} variant
seeds candidates on a Cartesian grid whose spacing scales with student
density, builds a greedy tour around each seed, and selects routes
minimizing total distance subject to full coverage, i.e.\ it minimizes
distance directly. Both are solved by branch-and-cut.

Formally, per-child builds $\mathcal{R}=\{R_1,\dots,R_n\}$ with
$R_i=\{i\}\cup\mathrm{NN}_{C-1}(i)$, the $C-1$ nearest neighbours of
$i$ under $D^{\mathrm{sym}}$, giving $|\mathcal{R}|=n$ candidates in
$O(n^2)$ time, and minimizes $\sum_j x_j$. Grid overlays a Cartesian
lattice whose spacing shrinks with student density, and for each
occupied cell greedily grows one capacity-feasible route around the
cell seed; its objective is $\sum_j \mathrm{cost}_j x_j$ with
$\mathrm{cost}_j$ the length of a greedy tour over $R_j$.

Because per-child never receives a distance term in its objective, its
routes are not a distance-minimizing candidate that happened to lose:
they are the cheapest-in-fleet-size candidate the solver was asked for.
Reading its longer routes as a quality gap against grid therefore
requires care. At least three factors can contribute, and our design
does not cleanly separate them:
\begin{itemize}
\item \textbf{Coverage vs.\ partitioning.} Eq.~(\ref{eq:setcover}) is a
\emph{covering} ($\ge1$) constraint, not a \emph{partitioning} ($=1$)
one, and the selected routes run as generated, so a student assigned to
two selected candidates is visited twice and both visits are charged:
at $N=400$ per-child selects $26.9$ capacity-filled routes, $538$ seats
for $400$ students, so roughly a quarter of visits are redundant.
\item \textbf{Candidate-pool size.} $\mathcal{R}$ is a heuristically
seeded subset of all $\binom{n}{\le C}$ feasible routes rather than an
exhaustive enumeration, so the branch-and-cut speed we measure reflects
a small pool, not an easy integer program.
\item \textbf{Objective misalignment.} Per-child optimizes fleet size
while grid optimizes distance directly, so comparing the two on the
distance column (Table~\ref{tab:fleet}) conflates pool quality with
objective choice. Isolating pool construction alone would require
running the fleet-size objective on the grid pool, or the distance
objective on the per-child pool, which we leave for future work and do
not attempt here.
\end{itemize}

\subsection{Adaptive large-neighbourhood search}\label{sec:alns}
\index{adaptive large-neighbourhood search (ALNS)}
The five strategies above cover cluster-first, route-second pipelines
and set-covering integer programs; ALNS adds a locality-aware
metaheuristic behind the same \texttt{Solver} interface
(Section~\ref{sec:framework}) as the sixth strategy in this benchmark,
so that a population-based, adaptive search is measured on the same
scenarios, metrics and solver interface as the other five.

ALNS starts from a feasible solution (the k-medoids baseline of
Section~\ref{sec:framework}) and repeatedly \emph{destroys} and
\emph{repairs} part of it (Algorithm~\ref{alg:alns}). A destroy
operator removes a subset of students from their current routes, drawn
from a pool that includes random removal, worst-cost removal (students
whose removal reduces route length the most) and Shaw-style related
removal (students that are close in $D^{\mathrm{sym}}$ or on the same
route, so removing them together opens a genuinely new neighbourhood).
A repair operator then reinserts the removed students, by cheapest
feasible insertion or a regret-$k$ criterion that anticipates future
insertion cost, respecting the capacity constraint of
Eq.~(\ref{eq:assign-cap}) at every step. Each iteration selects one
destroy and one repair operator, evaluates the resulting solution
against $L_{\mathrm{total}}$ (Eq.~\ref{eq:obj}), and accepts or rejects
it under a simulated-annealing-style criterion so that the search can
escape local optima without discarding the structure that
cluster-first, route-second exploits. Operator weights are updated
adaptively from their recent success, which is the property that gives
the method its name and is what should let it recover the locality
that a topology-agnostic representation, such as the flat-permutation
genetic encoding an earlier iteration of this benchmark tested and
discarded for scaling poorly, lacks.

\begin{algorithm}[t]
\caption{Adaptive large-neighbourhood search (ALNS)}
\label{alg:alns}
\footnotesize
\begin{algorithmic}[1]
\Require initial solution $r_0$ (k-medoids partition, NN$+$2-opt tours),
destroy operators $\Omega^-$, repair operators $\Omega^+$, time budget
$B_{\mathrm{A}}=k\cdot B$, cooling schedule $T_t=T_0\gamma^{\,t}$
\State $r \gets r_0$; \quad $r^{\ast}\gets r_0$
\State initialize operator weights $w_\omega \gets 1$ for all
$\omega\in\Omega^-\cup\Omega^+$
\For{$t = 1,2,\dots$ \textbf{while} elapsed $< B_{\mathrm{A}}$}
  \State select $\omega^-\in\Omega^-$, $\omega^+\in\Omega^+$
         with probability $\propto w_{\omega^-},w_{\omega^+}$
  \State $r' \gets \omega^+(\omega^-(r))$
         \Comment{destroy a subset of students, then repair}
  \If{$L_{\mathrm{total}}(r') < L_{\mathrm{total}}(r^{\ast})$}
     \State $r^{\ast}\gets r'$
  \EndIf
  \If{$L_{\mathrm{total}}(r')\le L_{\mathrm{total}}(r)$ \textbf{or}
      $\mathrm{rand}()<\exp\!\big(-[L_{\mathrm{total}}(r')-L_{\mathrm{total}}(r)]/T_t\big)$}
     \State $r \gets r'$ \Comment{simulated-annealing acceptance}
  \EndIf
  \State update $w_{\omega^-},w_{\omega^+}$ from the score of this
         iteration \Comment{adaptive weights}
\EndFor
\State \Return $r^{\ast}$
\end{algorithmic}
\end{algorithm}

\paragraph{Parameters and budget.}
The search is \emph{budget-matched} to the clustering pipelines: an
instance gets $B_{\mathrm{A}}=k\cdot B$ seconds of wall clock with
$B=5$~s, exactly the per-cluster TSP budget K-medoids and Sectorial spend
inside OR-Tools, so $15$, $25$, $50$ and $100$~s at
$N=50,100,200,400$. Any difference in solution quality is therefore
attributable to the search rather than to compute. Because the budget is
wall clock, the initial solution must be cheap: it reuses the baseline's
capacitated k-medoids partition
(Eqs.~(\ref{eq:assign-obj})--(\ref{eq:assign-cap})) but routes each
cluster with a nearest-neighbour tour plus 2-opt instead of calling
OR-Tools, costing $0.01$--$1.11$~s from $N=50$ to $N=400$ ($\le1.1\%$ of
the budget). Within it the search completes $15{,}095$, $13{,}374$,
$11{,}485$ and $10{,}353$ iterations respectively.

Each iteration removes $q\sim\mathcal{U}[0.10n,\,0.25n]$ students,
capped at $60$. The worst- and Shaw-removal operators pick from their
ranked candidate lists with the randomisation
$L[\lfloor |L|\cdot y^{p}\rfloor]$, $y\sim\mathcal{U}(0,1)$ and
$p=3$, so the choice is greedy but not deterministic; repair uses regret-2.
Operator weights start at $w_\omega=1$ and are updated as the convex
combination $w_\omega \gets \theta w_\omega + (1-\theta)s_j$ with
reaction factor $\theta=0.8$ and scores
$(s_1,\dots,s_4)=(5,2,1,0.5)$ for a new global best, an improvement on the
current solution, an accepted solution and a rejected one. Acceptance is
simulated annealing with geometric cooling $T_t=T_0\gamma^{\,t}$, where
$T_0$ is fitted so that a solution $5\%$ worse than $r_0$ is accepted with
probability $0.5$ at the start, and $\gamma$ anneals to $T=1$ over a
$25{,}000$-iteration horizon. The search stops when the budget is
exhausted and returns the best solution found. All of these values are
fixed across the entire sweep, with no per-instance tuning, exactly as for
the other five strategies (Sec.~\ref{sec:setup}).

\subsection{Complexity analysis}\label{sec:complexity}
Table~\ref{tab:complexity} summarizes each strategy's asymptotic
worst-case cost alongside the stage that actually dominates its measured
runtime, where $n$ is the number of students, $k=\lceil n/C\rceil$ the
number of buses and $T$ the iteration cap. The distinction matters,
because for the two flow-based pipelines the asymptotically expensive
stage is \emph{not} the one that governs the wall-clock times in
Section~\ref{sec:results}.

We timed the stages rather than assuming which dominates. For K-medoids
and Sectorial the minimum-cost-flow assignment takes $0.04\%$ of solver
time at $N=50$ and only $1.1\%$ at $N=400$ ($1.09$ of $101.13$~s), with
FasterPAM under $0.02$~s; the per-cluster TSP takes the remaining
$98.9$--$99.96\%$ and is \emph{budget-bound}, costing almost exactly
$k\cdot B$ ($15.03$, $25.01$, $50.01$, $100.03$~s for $k=3,5,10,20$ at
$B=5$~s). As $k=\lceil n/C\rceil$, latency grows \emph{linearly} in $n$,
as Table~\ref{tab:latency} shows. The $O(T\,n^3)$ flow bound is real but
latent: its share rises two orders of magnitude from $N=50$ to
$N=400$, so it dominates only far beyond the sizes studied.

ALNS is budget-bound in the same way but for a different reason: the time
limit is imposed directly rather than inherited from a sub-solver, so the
question its timers answer is not \emph{which stage dominates} but
\emph{how many iterations the budget buys}. Repair is the bottleneck at
every density, taking $67.0$, $71.2$, $71.4$ and $65.4\%$ of the run at
$N=50,100,200,400$ against $17.0$--$28.9\%$ for destroy; construction of
the initial solution never exceeds $1.1\%$, and the remaining
$4.6$--$15.9\%$ covers acceptance, operator-weight updates and objective
evaluation. Repair dominates because each insertion is scored over every
route and every position, an $O(k\,|r|)$ sweep per removed student, whereas
the removal operators need one pass over the current routes. Since the
budget grows as $k\cdot B$ while per-iteration cost grows with $n$, the
realised iteration count falls only from $15{,}095$ to $10{,}353$ across a
sweep in which $n$ grows eightfold.

\begin{table}[t]
\centering
\caption{Asymptotic worst-case cost per strategy, and the stage that
dominates the \emph{measured} runtime, from per-stage timers at $N{=}400$
($B$ = per-cluster TSP budget, $5$~s; $B_z$ = per-zone CVRPTW budget,
$20$~s). $^\dagger$~inferred from latency being invariant in $n$, not timed.}
\label{tab:complexity}
\small
\setlength{\tabcolsep}{4pt}
\begin{tabular}{lll}
\toprule
\textbf{Strategy} & \textbf{Asymptotic worst case} &
\textbf{Measured dominant stage}\\
\midrule
K-medoids + TSP & $O(T\,n^3)$, min-cost flow & TSP budget $k\,B$ ($98.9\%$)\\
Sectorial       & $O(T\,n^3)$, sweep $+$ flow & TSP budget $k\,B$ ($99.0\%$)\\
CVRPTW (zonal)  & $O(n^2)$--$O(n^3)$, GLS & zone budget $|Z|\,B_z^{\;\dagger}$\\
Set Cover (per-child) & ILP, exp.; $O(n^2)$ pool & ILP $67\%$, pool $28\%$\\
Set Cover (grid)      & ILP, exp.; $O(n)$ seeds & ILP ($>98\%$)\\
ALNS & $O(I\,n)$--$O(I\,n\log n)$, operator-dependent & repair $65\%$, destroy $29\%$\\
\bottomrule
\end{tabular}
\end{table}

\subsection{Experimental setup}\label{sec:setup}
Scenarios are generated on the real drivable road network of
Medell\'in and the Valle de Aburr\'a, extracted from OpenStreetMap with
OSMnx~\cite{Boeing2017OSMnx} and restricted to the largest strongly
connected component. The depot is fixed at a university campus, and
$N$ student stops are sampled uniformly from graph nodes lying inside
the union of the six service-zone polygons, so that all instances stay
within the actual metropolitan footprint. Directed shortest-path
distance and travel-time matrices are precomputed with Dijkstra's
algorithm. The bus capacity is $C=20$.

We sweep four student densities $N\in\{50,100,200,400\}$ with eleven
seeds each, giving $44$ scenarios and $264$ runs ($6$ strategies
$\times\,44$). For every run we record: total fleet distance, wall-clock
latency, number of buses used, coverage (fraction of students served),
the silhouette coefficient~\cite{Rousseeuw1987Silhouettes} of the
clustering on the symmetrized distances, capacity violations, and the
longest single-route distance. To keep latency comparable, the
per-cluster TSP search is capped at $5$~s and the per-zone CVRPTW
search at $20$~s.

All strategies run at fixed settings across the sweep, with no
per-instance tuning, so the comparison reflects out-of-the-box
behaviour. K-medoids uses iteration cap $T=50$; Sectorial uses
$\alpha=0.6$, $\beta=0.4$, min--max radial normalization and $T=5$;
the zonal CVRPTW uses $75$~min maximum duration, $90$~s boarding and a
span-balancing coefficient of $300$; ALNS uses a destroy fraction of
$0.10$--$0.25$, reaction factor $\theta=0.8$ with scores $(5,2,1,0.5)$,
regret-2 repair and simulated-annealing acceptance fitted as described in
Section~\ref{sec:alns}, under a wall-clock budget of $k\cdot5$~s that
matches the clustering pipelines. These settings are held fixed across the
whole sweep, matching the no-per-instance-tuning protocol used for the
other five strategies.
Statistical significance across
strategies is assessed with the non-parametric Friedman test (blocks
$=$ the $44$ scenarios), and pairwise differences with the Wilcoxon
signed-rank test.

\subsubsection{External calibration.}\label{sec:external}
A benchmark of six in-house implementations on self-generated instances
establishes relative standings but says nothing about absolute quality. To
supply that reference point we run the ALNS solver, unchanged and through
the same \texttt{Solver} interface, on the $27$ instances of the Augerat
set~A~\cite{Augerat1995CVRP}, a classical CVRP family with published
optimal values, $n=31$ to $79$ customers and capacity $100$. The loader
reproduces the TSPLIB \texttt{EUC\_2D} convention exactly, that is,
integer distances $\mathrm{nint}(\sqrt{\Delta x^2+\Delta y^2})$, which we
verified by recomputing the cost of each published optimal solution under
our own matrix and recovering the published value for all $27$ instances.
Two differences from the SBRP setting are worth naming: these instances
carry non-unit customer demands $q_i$, so capacity is a demand sum rather
than a stop count, and each has a declared vehicle count that matches
$\lceil \sum_i q_i / C\rceil$ in every case, so ALNS optimizes over the
same fleet size the published optimum uses. Each instance gets a fixed
$60$~s budget. Only ALNS is run here: the clustering pipelines assume unit
demand, under which their capacitated assignment is an exact minimum-cost
flow, and generalising them to arbitrary demands would change the
algorithms being compared.

% =====================================================================
\section{Results}\label{sec:results}
% =====================================================================
Table~\ref{tab:fleet} reports fleet distance and latency, and
Table~\ref{tab:quality} the remaining quality metrics at the most
demanding density ($N=400$). Figures~\ref{fig:scalability}
and~\ref{fig:pareto} show how latency scales and the distance--latency
trade-off; Figure~\ref{fig:grid} contrasts the solution paradigms on a
shared scenario. All three are generated from the same $308$-run output
as the tables, so figure and table can not drift apart.

\begin{table}[t]
\centering
\caption{Total fleet distance (km, mean $\pm$ std over 11 seeds) and
latency (s, mean). Bold: distance minimum among full-coverage methods.}
\label{tab:fleet}
\label{tab:latency}
\setlength{\tabcolsep}{2.2pt}
\footnotesize
\begin{tabular}{l@{\hskip 4pt}cccc@{\hskip 7pt}cccc}
\toprule
& \multicolumn{4}{c@{\hskip 10pt}}{\textbf{Fleet distance (km)}}
& \multicolumn{4}{c}{\textbf{Latency (s)}}\\
\cmidrule(r){2-5}\cmidrule{6-9}
\textbf{Strategy} & $50$ & $100$ & $200$ & $400$
                  & $50$ & $100$ & $200$ & $400$\\
\midrule
K-medoids + TSP      & 105$\pm$5  & 163$\pm$18 & 271$\pm$21 & 460$\pm$26
  & 15.0 & 25.0 & 50.2 & 101.0\\
Sectorial            & 105$\pm$5  & 165$\pm$19 & 272$\pm$21 & 461$\pm$25
  & 15.0 & 25.0 & 50.2 & 101.0\\
CVRPTW (zonal)$^\ast$ & 125$\pm$7  & 182$\pm$12 & 286$\pm$16 & 471$\pm$13
  & 100.0 & 100.0 & 100.0 & 100.0\\
Set Cover (per-child) & 160$\pm$27 & 254$\pm$36 & 389$\pm$39 & 670$\pm$20
  & 0.02 & 0.02 & 0.04 & 0.16\\
Set Cover (grid)      & 134$\pm$8  & 217$\pm$22 & 349$\pm$26 & 594$\pm$33
  & 0.27 & 0.74 & 3.19 & 66.6\\
ALNS                 & \textbf{98$\pm$5}  & \textbf{155$\pm$19} & \textbf{251$\pm$19} & \textbf{428$\pm$24}
  & 15.0 & 25.0 & 50.0 & 100.0\\
\bottomrule
\end{tabular}
\\[3pt]
{\footnotesize $^\ast$ CVRPTW may leave ${\sim}1\%$ of students unserved
when the duration limit binds, so its distance is not strictly comparable
to the full-coverage methods.}
\end{table}

\begin{table}[t]
\centering
\caption{Quality metrics at $N{=}400$, mean over 11 seeds.}
\label{tab:quality}
\small
\begin{tabular}{lcccc}
\toprule
\textbf{Strategy} & \textbf{Buses} & \textbf{Coverage} & \textbf{Silhouette} & \textbf{Max route (km)}\\
\midrule
K-medoids + TSP      & 20.0 & 1.000 & \textbf{0.276} & 46.2\\
Sectorial            & 20.0 & 1.000 & 0.266 & 45.2\\
CVRPTW (zonal)       & 22.4 & 0.998 & 0.272 & \textbf{30.3}\\
Set Cover (per-child) & 26.9 & 1.000 & 0.230 & 50.8\\
Set Cover (grid)      & 23.2 & 1.000 & 0.199 & 56.6\\
ALNS                  & 20.0 & 1.000 & 0.068 & 43.3\\
\bottomrule
\end{tabular}
\end{table}

\paragraph{Statistical significance.}
The Friedman test rejects the null hypothesis of equal performance
across the six strategies for both fleet distance
($\chi^2=201.74$, $\mathrm{df}=5$, $p<10^{-40}$) and latency
($\chi^2=187.57$, $\mathrm{df}=5$, $p<10^{-37}$). On fleet distance the
mean ranks (1 = best) are $1.02$ for ALNS, $2.63$ for Sectorial and
$2.67$ for K-medoids, $3.68$ for CVRPTW, $5.05$ for Set Cover (grid) and
$5.95$ for Set Cover (per-child). ALNS's rank of $1.02$ reflects an almost
unbroken sweep: it returns the shortest fleet distance in $44$ of the $44$
scenarios against both clustering pipelines, and in $43$ of $44$ against
the zonal CVRPTW, whose single win comes on an instance where it leaves
students unserved.
The two clustering leaders remain statistically indistinguishable from
each other: a Wilcoxon signed-rank test between K-medoids and Sectorial
over the $44$ scenarios yields $p=0.36$, unaffected by which other
strategies are included in the comparison. On latency, Set Cover
(per-child) remains the fastest strategy in every scenario (mean rank
$1.00$), followed by Set Cover (grid) ($2.18$) and then ALNS ($2.98$),
ahead of K-medoids ($4.59$), Sectorial ($4.82$) and the zonal CVRPTW
($5.43$). ALNS outranking the clustering pipelines on latency is an
artefact worth stating plainly rather than a speed advantage: all three
are given the same $k\cdot5$~s, but ALNS stops exactly on its budget
whereas the OR-Tools calls overshoot theirs by about $1\%$
(Table~\ref{tab:fleet}). The quality/speed tension visible in
Figure~\ref{fig:pareto} therefore now has a strategy at the favourable
corner of both axes among the budgeted methods.

Since a Friedman rejection only says \emph{some} strategies differ, we
add the Nemenyi post-hoc~\cite{Demsar2006Statistical}, which separates
two strategies when their mean ranks differ by more than
$\mathrm{CD}=q_\alpha\sqrt{k(k+1)/6N}$. With $k=6$ strategies, $N=44$
blocks and $\alpha=0.05$ ($q_\alpha=2.850$ for $k=6$, the tabulated
Nemenyi critical value), $\mathrm{CD}=2.850\sqrt{6\cdot7/(6\cdot44)}
\approx 1.14$. There are $\binom{6}{2}=15$ pairwise comparisons in
total, of which $11$ clear $\mathrm{CD}$ on fleet distance. ALNS
separates from every other strategy: from Set Cover (per-child)
($\Delta=4.93$), Set Cover (grid) ($4.02$), the zonal CVRPTW ($2.66$) and,
decisively for the headline claim, from K-medoids ($1.65$) and Sectorial
($1.60$), both comfortably above $\mathrm{CD}=1.14$. The two clustering
pipelines in turn separate from Set Cover (per-child) ($3.28$--$3.33$) and
Set Cover (grid) ($2.38$--$2.42$), and the CVRPTW from both set-covering
variants ($1.36$ and $2.27$). Four pairs remain unresolved: K-medoids
versus Sectorial ($0.045$), each of them against the CVRPTW
($1.01$--$1.06$), and the two set-covering variants against each other
($0.91$). The design therefore resolves the comparison that matters most
here, ALNS against the previous leaders, while leaving the same
near-ties open that eleven seeds could not separate before.

\paragraph{Absolute quality against published optima.}
On the external instances of Section~\ref{sec:external} the same ALNS
configuration solves $19$ of $27$ to proven optimality and averages
$0.221\%$ above the optimum, with a median gap of exactly $0$ and a worst
case of $1.71\%$ (A-n64-k9); every returned solution is capacity-feasible,
covers all customers and uses the declared number of vehicles
(Table~\ref{tab:cvrplib}). The eight non-zero gaps are concentrated in the
larger instances, $n\ge52$, which is where a $60$~s budget buys
proportionally fewer iterations. This does not make our ALNS
state-of-the-art, dedicated CVRP solvers close set~A routinely, but it does
bound what the rest of this paper leaves open: the strategy that wins our
benchmark is itself within a fraction of a percent of optimal on
independent instances, so the margins reported in
Table~\ref{tab:fleet} are differences between a near-optimal method and
the rest, not between six equally unvalidated ones.

\begin{table}[t]
\centering
\caption{ALNS against published optima on Augerat set~A, $60$~s per
instance. Gap is $100(\mathrm{ALNS}-\mathrm{opt})/\mathrm{opt}$.
All $27$ solutions are feasible and use the declared vehicle count.}
\label{tab:cvrplib}
\scriptsize
\setlength{\tabcolsep}{2.5pt}
\begin{tabular}{lrrr@{\hskip 8pt}lrrr@{\hskip 8pt}lrrr}
\toprule
\textbf{Inst.} & $n$ & \textbf{opt} & \textbf{\%} &
\textbf{Inst.} & $n$ & \textbf{opt} & \textbf{\%} &
\textbf{Inst.} & $n$ & \textbf{opt} & \textbf{\%}\\
\midrule
A-n32-k5 & 31 & 784 & 0.00 & A-n33-k5 & 32 & 661 & 0.00 & A-n33-k6 & 32 & 742 & 0.00 \\
A-n34-k5 & 33 & 778 & 0.00 & A-n36-k5 & 35 & 799 & 0.00 & A-n37-k5 & 36 & 669 & 0.00 \\
A-n37-k6 & 36 & 949 & 0.00 & A-n38-k5 & 37 & 730 & 0.00 & A-n39-k5 & 38 & 822 & 0.00 \\
A-n39-k6 & 38 & 831 & 0.00 & A-n44-k6 & 43 & 937 & 0.00 & A-n45-k6 & 44 & 944 & 0.00 \\
A-n45-k7 & 44 & 1146 & 0.00 & A-n46-k7 & 45 & 914 & 0.00 & A-n48-k7 & 47 & 1073 & 0.00 \\
A-n53-k7 & 52 & 1010 & 0.69 & A-n54-k7 & 53 & 1167 & 0.00 & A-n55-k9 & 54 & 1073 & 0.00 \\
A-n60-k9 & 59 & 1354 & 0.66 & A-n61-k9 & 60 & 1034 & 0.10 & A-n62-k8 & 61 & 1288 & 0.00 \\
A-n63-k9 & 62 & 1616 & 0.68 & A-n63-k10 & 62 & 1314 & 0.38 & A-n64-k9 & 63 & 1401 & 1.71 \\
A-n65-k9 & 64 & 1174 & 0.00 & A-n69-k9 & 68 & 1159 & 0.78 & A-n80-k10 & 79 & 1763 & 0.96 \\
\midrule
\multicolumn{12}{l}{\textbf{19/27 optimal} \quad mean $0.221\%$ \quad
median $0.000\%$ \quad max $1.71\%$}\\
\bottomrule
\end{tabular}
\end{table}

\begin{figure}[t]
\centering
\begin{minipage}[t]{0.49\textwidth}
  \centering
  \includegraphics[width=\textwidth]{figures/scalability.pdf}
  \Description{Log-log line plot of mean solver latency in seconds
  against the number of students, one line per strategy.}
  \caption{Mean latency vs.\ $N$ (log--log). Clustering pipelines grow
  linearly in the number of buses because each is routed under a fixed
  per-cluster TSP budget; the flow assignment contributes at most
  $1.1\%$ of their runtime. K-medoids and Sectorial have near-identical
  latencies, so Sectorial is dashed with hollow markers over the solid
  K-medoids curve, and both series are present throughout.}
  \label{fig:scalability}
\end{minipage}\hfill
\begin{minipage}[t]{0.49\textwidth}
  \centering
  \includegraphics[width=\textwidth]{figures/pareto_N400.pdf}
  \Description{Scatter plot of fleet distance against latency at
  N=400, one marker per strategy, with the two leading strategies
  nearly overlapping in the lower-left corner.}
  \caption{Distance--latency trade-off at $N{=}400$ (lower-left is
  better). The two leaders are all but coincident, differing by
  $1.0$~km ($0.22\%$) and $0.03$~s, so Sectorial is drawn as a hollow
  ring over the filled K-medoids marker rather than hidden beneath it,
  and the annotation reports the measured gap.}
  \label{fig:pareto}
\end{minipage}
\end{figure}

\begin{figure}[t]
\centering
\includegraphics[width=0.95\textwidth]{figures/comparison_row_N200.png}
\Description{Side-by-side maps of the same scenario: wedge-shaped
sectorial clusters, overlapping grid set-covering routes, and the
ALNS solution, drawn over the real road network.}
\caption{Three representatives of the solution paradigms compared here
on a shared scenario ($N{=}200$, seed 1). Left: Sectorial produces
compact, wedge-shaped clusters radiating from the depot. Centre: the
grid set-covering ILP fills every selected route to capacity, so its
routes overlap and some students are visited twice. Right: ALNS serves the
same ten buses as Sectorial but abandons the wedge structure; its routes
are elongated corridors that follow the valley and interleave near the
depot, so neighbouring buses share territory wherever that shortens the
tours. This interleaving is exactly what the low silhouette of
Table~\ref{tab:quality} measures, and it coincides with the shortest
fleet distance rather than the longest.}
\label{fig:grid}
\end{figure}

% =====================================================================
\section{Discussion}\label{sec:discussion}
% =====================================================================
\paragraph{Clustering is a good heuristic, not the objective.}
Among the cluster-first, route-second pipelines the story is the familiar
one: K-medoids and Sectorial produce spatially compact, well-separated
clusters, K-medoids attains the highest silhouette at $N=400$ ($0.276$;
Table~\ref{tab:quality}), and compact clusters shorten the inter-stop
distances the per-cluster TSP must traverse, so a good partition does most
of the optimization before routing begins~\cite{CakirRevisitingProblem,Comert2018AStudy}.
The cheap angular score (Eq.~\ref{eq:sectorial-score}) recovers
essentially the quality of medoid optimization (Wilcoxon $p=0.36$); the
$1.0$~km that separates them at $N{=}400$ is well inside seed noise.

ALNS breaks that story, and this is the most consequential finding of the
present study. It attains the \emph{shortest} fleet distance at every
density while producing by far the \emph{least} cohesive partition:
silhouette $0.068$ at $N=400$, a quarter of K-medoids' $0.276$, and
falling monotonically with scale ($0.325$, $0.265$, $0.174$, $0.068$ from
$N=50$ to $N=400$) even as its distance advantage holds at
$4.8$--$7.3\%$. Cohesion and cost are therefore not two views of the same
quantity. A capacitated partition chosen to be compact is a good
\emph{heuristic} for short routes, because it makes each TSP easy; but it
is a constraint on the solution, and a search that optimizes
$L_{\mathrm{total}}$ directly will spend that compactness whenever
interleaving two buses' territories buys a shorter tour
(Fig.~\ref{fig:grid}, right).

This also revises how the correlation between cohesion and cost should be
read. Measured over the runs within a density, silhouette and fleet
distance correlate at $r=-0.95$ at $N{=}400$ when the sixth strategy is a
flat-permutation genetic algorithm, as in an earlier iteration of this
benchmark. That coefficient does not survive substitution: with ALNS in
its place the same computation gives $r=+0.09$ ($p=0.48$), and no density
shows a significant negative association ($r=-0.35$ to $+0.09$). The
strong correlation was carried by a single strategy that was simultaneously
incoherent and long, not by a general law relating the two. We therefore
report cohesion as a descriptive property of a solution rather than as a
predictor of its cost, and we would caution against selecting an SBRP
solver on a clustering-quality score.

\paragraph{The CVRPTW trade-off.}
The zonal CVRPTW yields competitive distances but is not free: it
leaves up to ${\sim}1\%$ of students unserved when the $75$-minute
duration limit binds, and it uses more buses than the optimum
($22.4$ vs.\ $20$ at $N=400$). In exchange it produces by far the most
\emph{balanced} routes: its longest route ($30.3$ km) is
$33$--$46\%$ shorter than any other method's longest route among the
five strategies compared here (Table~\ref{tab:quality}), because the
span-balancing penalty and the explicit time dimension cap individual
ride times. Where bounded student ride time is a hard requirement, this
is the only strategy that delivers it directly.

\paragraph{Set covering: a speed/quality trade-off with two different
objectives.}
The set-cover programs are the speed champions: the per-child variant
solves every instance in under a quarter of a second, and the grid
variant is an order of magnitude slower but still far faster than the
clustering pipelines at $N\le200$. On distance they occupy the last two
places of the six, with mean ranks $5.05$ for grid and $5.95$ for
per-child; on latency they occupy the first two, $2.18$ and $1.00$. As
discussed
in Section~\ref{sec:setcover-two}, however,
the two variants are not a controlled test of the same design choice:
per-child minimizes fleet size and never receives a distance signal, so
its longer routes are an expected consequence of the objective it was
given, not solid evidence that its candidate pool is worse than grid's
for the same task. Grid trades some of per-child's speed for a direct
distance objective, but its own latency becomes erratic at $N=400$
($66.6\pm34.7$~s) as the candidate pool, and hence the integer program,
grows. Both remain the right choice when a feasible plan is needed in
near-real-time, e.g.\ for interactive what-if analysis, with the
caveat that comparing their route quality to each other should be read
as a comparison of two policies, not a controlled ablation of pool
construction.

\paragraph{Limitations.}
Eleven seeds per density give $44$ blocks. With six strategies there
are $\binom{6}{2}=15$ pairwise comparisons in total, of which the
Nemenyi critical difference separates $11$ on fleet distance
(Section~\ref{sec:results}). Four remain unresolved: K-medoids versus
Sectorial ($\Delta=0.045$), each of those two against the zonal CVRPTW
($1.01$ and $1.06$, both just under $\mathrm{CD}=1.14$), and the two
set-covering variants against each other ($0.91$). The K-medoids versus
Sectorial tie is a genuine equivalence by every test we applied; the other
three are power limits of a $44$-block design, not evidence of
equivalence, and we do not report them as such. Every comparison involving
ALNS is resolved.

The external reference point is now partial rather than absent. All six
strategies are our own implementations on instances we generate, and
Section~\ref{sec:external} calibrates only \emph{one} of them, ALNS,
and only against published \emph{CVRP} optima, not SBRP instances. Three
gaps remain. The other five strategies are still uncalibrated in absolute
terms, so their standings rest on the comparison with ALNS rather than on
any external value. Set~A instances are small ($n\le79$) and Euclidean,
whereas our scenarios are larger and run on a directed road network, so
the near-optimality shown there need not carry to $N=400$ in Medell\'in.
And we compare against published optimal \emph{values}, not against other
authors' running implementations, so we can state how far our solutions
are from optimal but not how our runtimes compare with theirs. A related, narrower limitation is that our
results speak to \emph{one} ALNS configuration (Section~\ref{sec:alns}):
a different operator mix, acceptance schedule or iteration budget could
shift its standing, so its ranking here should not be read as a bound
on what ALNS, or the metaheuristic family more broadly, can achieve on
this problem. Finally, the
benchmark models a single school with static demand and a homogeneous
fleet, and all strategies share one OR-Tools budget, so absolute
latencies reflect that budget rather than an intrinsic bound.

% =====================================================================
\section{Conclusions}\label{sec:conclusions}
% =====================================================================
We presented a unified, road-network-grounded benchmark of six SBRP
solution strategies evaluated on identical scenarios under a common
solver interface, with a formal model per strategy, complexity
validated against measured timings, and non-parametric statistical
testing. The central finding is that capacitated k-medoids and a
multi-objective sectorial clustering are tied for the shortest fleet
distance among the five clustering- and covering-based strategies with
established results, because both build compact clusters that make the
downstream TSP easy. Set-covering integer programs occupy the fast end
of the speed/quality trade-off, at different points of it because they
optimize different objectives (fleet size vs.\ distance), a distinction
we make explicit so it is not mistaken for a simple quality gap. A
zonal CVRPTW is the method of choice when bounded ride time is
mandatory. A locality-aware metaheuristic, an adaptive
large-neighbourhood search (ALNS, Section~\ref{sec:alns}), is included
as the sixth strategy, tested under the same harness and metrics used
throughout this paper. Given the same wall clock as the clustering
pipelines it is the strongest strategy we measured, taking the shortest
fleet distance in all $44$ scenarios, $4.8$--$7.3\%$ below the previous
leaders, and separating from them by more than the Nemenyi critical
difference; calibration against $27$ published CVRP instances puts it a
mean $0.22\%$ from the known optima, with $19$ solved to proven
optimality, which is the external reference point the rest of this
benchmark lacks. Its partitions are also the least cohesive of the six, so
the practical lesson is not that clustering was the wrong idea but that
cluster compactness is a means to short routes and not a proxy for
them.

Future work will (i)~extend the external calibration of
Section~\ref{sec:external} to the remaining five strategies, to larger
and road-network-based instances, and to running comparisons against other
authors' implementations rather than published values alone; (ii)~broaden the statistical design with
more seeds to resolve the pairs the Nemenyi test leaves open;
(iii)~explore the sensitivity of the ALNS results to its operator mix,
acceptance schedule and iteration budget (Section~\ref{sec:alns}),
since the present benchmark reports a single configuration rather than
a tuned one; (iv)~extend the harness to multi-school and
dynamic-demand variants; and (v)~convert the vehicle-kilometre results
into emissions and operating cost using fleet-specific factors, so that
solver selection can be argued directly in the sustainability terms
that motivate it.

\subsubsection{Use of generative AI.}
Generative AI assisted with drafting and copy-editing the text, with the
scripts that compute the reported statistics and render the figures from the
benchmark output, and with literature search. It did not generate experimental
results: every reported number is the output of the described code run on the
described instances, and the authors verified all numbers, figures and
references against the underlying data and primary sources. The authors assume
full responsibility for the content of this paper.

\ethics{Competing Interests}{\todo{Confirmar: los autores declaran no
tener conflictos de interes relevantes para el contenido de este
capitulo, o completar segun corresponda.}}

\ethics{Ethics Approval}{This study does not involve human participants
or animal subjects: all student locations are synthetic scenarios
sampled from publicly available OpenStreetMap road-network data, not
records of real students. \todo{Confirmar si la universidad requiere una
declaracion adicional.}}

% La plantilla oficial de camera-ready (CITIS-Final-Paper.tex) no usa
% \bibliographystyle/\bibliography: incluye la bibliografia manualmente
% via %%%%%%%%%%%%%%%%%%%%%%%% References.tex %%%%%%%%%%%%%%%%%%%%%%%%%%%%%%
% References for the paper:
% Optimization of School Bus Routing Strategies in Medellin
%
% Springer-style numbered references.
% All entries below correspond to citations used in optimization.tex.
%%%%%%%%%%%%%%%%%%%%%%%%%%%%%%%%%%%%%%%%%%%%%%%%%%%%%%%%%%%%%%%%%%%%%%%

\begin{thebibliography}{99}%

\bibitem{Park2010TheReview}
Park, J., Kim, B.-I.: The school bus routing problem: A review. Eur. J. Oper. Res. \textbf{202}(2), 311--319 (2010). doi: 10.1016/j.ejor.2009.05.017

\bibitem{Augerat1995CVRP}
Augerat, P., Belenguer, J.M., Benavent, E., Corber\'an, A., Naddef, D., Rinaldi, G.: Computational results with a branch-and-cut code for the capacitated vehicle routing problem. Research Report RR 949-M, Universit\'e Joseph Fourier, Grenoble (1995)

\bibitem{DantzigRamser1959Truck}
Dantzig, G.B., Ramser, J.H.: The truck dispatching problem. Manage. Sci. \textbf{6}(1), 80--91 (1959). doi: 10.1287/mnsc.6.1.80

\bibitem{Ellegood2020SchoolDirections}
Ellegood, W.A., Solomon, S., North, J., Campbell, J.F.: School bus routing problem: Contemporary trends and research directions. Omega \textbf{95}, 102056 (2020). doi: 10.1016/j.omega.2019.03.014

\bibitem{Croes1958TwoOpt}
Croes, G.A.: A method for solving traveling-salesman problems. Oper. Res. \textbf{6}(6), 791--812 (1958). doi: 10.1287/opre.6.6.791

\bibitem{HelsgaunAnHeuristic}
Helsgaun, K.: An effective implementation of the Lin--Kernighan traveling salesman heuristic. Eur. J. Oper. Res. \textbf{126}(1), 106--130 (2000). doi: 10.1016/S0377-2217(99)00284-2

\bibitem{CakirRevisitingProblem}
Cakir, F., Street, W.N., Thomas, B.W.: Revisiting cluster first, route second for the vehicle routing problem. Working paper (2015).

\bibitem{Comert2018AStudy}
Comert, S.E., Yazgan, H.R., K{\i}r, S., Yener, F.: A cluster first-route second approach for a capacitated vehicle routing problem: A case study. Int. J. Procurement Manag. \textbf{11}(4), 399--419 (2018). doi: 10.1504/IJPM.2018.092766

\bibitem{Corberan2002HeuristicObjectives}
Corberán, A., Fernández, E., Laguna, M., Martí, R.: Heuristic solutions to the problem of routing school buses with multiple objectives. J. Oper. Res. Soc. \textbf{53}(4), 427--435 (2002). doi: 10.1057/palgrave.jors.2601324

\bibitem{Baranwal2016DAVRPTW}
Baranwal, M., Parekh, P.M., Marla, L., Salapaka, S.M., Beck, C.L.: Vehicle routing problem with time windows: A deterministic annealing approach. In: 2016 American Control Conference (ACC), pp. 790--795. IEEE (2016). doi: 10.1109/ACC.2016.7525010

\bibitem{Hatamlou2018BlackHole}
Hatamlou, A.: Solving travelling salesman problem using black hole algorithm. Soft Comput. \textbf{22}(24), 8167--8175 (2018). doi: 10.1007/s00500-017-2760-y

\bibitem{Sabet2022GreenVRP}
Sabet, S., Farooq, B.: Green vehicle routing problem: State of the art and future directions. IEEE Access \textbf{10}, 101622--101642 (2022). doi: 10.1109/ACCESS.2022.3208899

\bibitem{Wang2025SBRPReview}
Wang, K., Yang, Y., Liao, M., Chen, A.: The school bus routing problem: A comprehensive review. Eur. Transp. Stud. \textbf{2}, 100043 (2025). doi: 10.1016/j.ets.2025.100043

\bibitem{KaufmanRousseeuw1990PAM}
Kaufman, L., Rousseeuw, P.J.: Finding Groups in Data: An Introduction to Cluster Analysis. Wiley, New York (1990)

\bibitem{Akitaya2021Subtrajectory}
Akitaya, H.A., Brüning, F., Chambers, E., Driemel, A.: Subtrajectory clustering: Finding set covers for set systems of subcurves. Comput. Geom. Topol. \textbf{2}(1), 7 (2023). doi: 10.57717/cgt.v2i1.7

\bibitem{ClarkeWright1964Savings}
Clarke, G., Wright, J.W.: Scheduling of vehicles from a central depot to a number of delivery points. Oper. Res. \textbf{12}(4), 568--581 (1964). doi: 10.1287/opre.12.4.568

\bibitem{PerronFurnon2023ORTools}
Didier, F., Perron, L., Mohajeri, S., Gay, S.A., Cuvelier, T., Furnon, V.: OR-Tools' vehicle routing solver: A generic constraint-programming solver with heuristic search for routing problems. In: Congrès de la Société française de recherche opérationnelle et d'aide à la décision (ROADEF 2023) (2023)

\bibitem{Voudouris1999GLS}
Voudouris, C., Tsang, E.: Guided local search and its application to the traveling salesman problem. Eur. J. Oper. Res. \textbf{113}(2), 469--499 (1999). doi: 10.1016/S0377-2217(98)00099-X

\bibitem{Holland1992Adaptation}
Holland, J.H.: Adaptation in Natural and Artificial Systems: An Introductory Analysis with Applications to Biology, Control, and Artificial Intelligence. MIT Press, Cambridge, MA (1992). doi: 10.7551/mitpress/1090.001.0001

\bibitem{Joshi2019GCN}
Joshi, C.K., Laurent, T., Bresson, X.: An efficient graph convolutional network technique for the travelling salesman problem. arXiv:1906.01227 (2019)

\bibitem{Dahiya2018TSPReview}
Dahiya, C., Sangwan, S.: Literature review on travelling salesman problem. Int. J. Res. \textbf{5}(16), 1152--1155 (2018)

\bibitem{Wu2024NCOSurvey}
Wu, X., Wang, D., Wen, L., Xiao, Y., Wu, C., Wu, Y., Yu, C., Maskell, D.L., Zhou, Y.: Neural combinatorial optimization algorithms for solving vehicle routing problems: A comprehensive survey with perspectives. arXiv:2406.00415 (2024). doi: 10.48550/arXiv.2406.00415

\bibitem{Boeing2017OSMnx}
Boeing, G.: OSMnx: New methods for acquiring, constructing, analyzing, and visualizing complex street networks. Comput. Environ. Urban Syst. \textbf{65}, 126--139 (2017). doi: 10.1016/j.compenvurbsys.2017.05.004

\bibitem{Pedregosa2011Sklearn}
Pedregosa, F., Varoquaux, G., Gramfort, A., et al.: Scikit-learn: Machine learning in Python. J. Mach. Learn. Res. \textbf{12}, 2825--2830 (2011)

\bibitem{Rousseeuw1987Silhouettes}
Rousseeuw, P.J.: Silhouettes: A graphical aid to the interpretation and validation of cluster analysis. J. Comput. Appl. Math. \textbf{20}, 53--65 (1987). doi: 10.1016/0377-0427(87)90125-7

\bibitem{Demsar2006Statistical}
Demšar, J.: Statistical comparisons of classifiers over multiple data sets. J. Mach. Learn. Res. \textbf{7}, 1--30 (2006)

\end{thebibliography}
, que ya contiene el entorno thebibliography.
% Mantener ambos mecanismos duplica la bibliografia y rompe la compilacion
% si se ejecuta bibtex/latexmk, ademas refs.bib no esta en esta carpeta.
% \bibliographystyle{spmpsci}
% \bibliography{refs}
% \todo{Springer recomienda, para la version camera-ready, correr
% bibtex localmente y pegar el .bbl resultante aqui en lugar de
% \bibliography{refs}, o subir refs.bib junto con el .bst
% correspondiente (ver instrucciones de LaTeX de Springer Nature).}
%%%%%%%%%%%%%%%%%%%%%%%% References.tex %%%%%%%%%%%%%%%%%%%%%%%%%%%%%%
% References for the paper:
% Optimization of School Bus Routing Strategies in Medellin
%
% Springer-style numbered references.
% All entries below correspond to citations used in optimization.tex.
%%%%%%%%%%%%%%%%%%%%%%%%%%%%%%%%%%%%%%%%%%%%%%%%%%%%%%%%%%%%%%%%%%%%%%%

\begin{thebibliography}{99}%

\bibitem{Park2010TheReview}
Park, J., Kim, B.-I.: The school bus routing problem: A review. Eur. J. Oper. Res. \textbf{202}(2), 311--319 (2010). doi: 10.1016/j.ejor.2009.05.017

\bibitem{Augerat1995CVRP}
Augerat, P., Belenguer, J.M., Benavent, E., Corber\'an, A., Naddef, D., Rinaldi, G.: Computational results with a branch-and-cut code for the capacitated vehicle routing problem. Research Report RR 949-M, Universit\'e Joseph Fourier, Grenoble (1995)

\bibitem{DantzigRamser1959Truck}
Dantzig, G.B., Ramser, J.H.: The truck dispatching problem. Manage. Sci. \textbf{6}(1), 80--91 (1959). doi: 10.1287/mnsc.6.1.80

\bibitem{Ellegood2020SchoolDirections}
Ellegood, W.A., Solomon, S., North, J., Campbell, J.F.: School bus routing problem: Contemporary trends and research directions. Omega \textbf{95}, 102056 (2020). doi: 10.1016/j.omega.2019.03.014

\bibitem{Croes1958TwoOpt}
Croes, G.A.: A method for solving traveling-salesman problems. Oper. Res. \textbf{6}(6), 791--812 (1958). doi: 10.1287/opre.6.6.791

\bibitem{HelsgaunAnHeuristic}
Helsgaun, K.: An effective implementation of the Lin--Kernighan traveling salesman heuristic. Eur. J. Oper. Res. \textbf{126}(1), 106--130 (2000). doi: 10.1016/S0377-2217(99)00284-2

\bibitem{CakirRevisitingProblem}
Cakir, F., Street, W.N., Thomas, B.W.: Revisiting cluster first, route second for the vehicle routing problem. Working paper (2015).

\bibitem{Comert2018AStudy}
Comert, S.E., Yazgan, H.R., K{\i}r, S., Yener, F.: A cluster first-route second approach for a capacitated vehicle routing problem: A case study. Int. J. Procurement Manag. \textbf{11}(4), 399--419 (2018). doi: 10.1504/IJPM.2018.092766

\bibitem{Corberan2002HeuristicObjectives}
Corberán, A., Fernández, E., Laguna, M., Martí, R.: Heuristic solutions to the problem of routing school buses with multiple objectives. J. Oper. Res. Soc. \textbf{53}(4), 427--435 (2002). doi: 10.1057/palgrave.jors.2601324

\bibitem{Baranwal2016DAVRPTW}
Baranwal, M., Parekh, P.M., Marla, L., Salapaka, S.M., Beck, C.L.: Vehicle routing problem with time windows: A deterministic annealing approach. In: 2016 American Control Conference (ACC), pp. 790--795. IEEE (2016). doi: 10.1109/ACC.2016.7525010

\bibitem{Hatamlou2018BlackHole}
Hatamlou, A.: Solving travelling salesman problem using black hole algorithm. Soft Comput. \textbf{22}(24), 8167--8175 (2018). doi: 10.1007/s00500-017-2760-y

\bibitem{Sabet2022GreenVRP}
Sabet, S., Farooq, B.: Green vehicle routing problem: State of the art and future directions. IEEE Access \textbf{10}, 101622--101642 (2022). doi: 10.1109/ACCESS.2022.3208899

\bibitem{Wang2025SBRPReview}
Wang, K., Yang, Y., Liao, M., Chen, A.: The school bus routing problem: A comprehensive review. Eur. Transp. Stud. \textbf{2}, 100043 (2025). doi: 10.1016/j.ets.2025.100043

\bibitem{KaufmanRousseeuw1990PAM}
Kaufman, L., Rousseeuw, P.J.: Finding Groups in Data: An Introduction to Cluster Analysis. Wiley, New York (1990)

\bibitem{Akitaya2021Subtrajectory}
Akitaya, H.A., Brüning, F., Chambers, E., Driemel, A.: Subtrajectory clustering: Finding set covers for set systems of subcurves. Comput. Geom. Topol. \textbf{2}(1), 7 (2023). doi: 10.57717/cgt.v2i1.7

\bibitem{ClarkeWright1964Savings}
Clarke, G., Wright, J.W.: Scheduling of vehicles from a central depot to a number of delivery points. Oper. Res. \textbf{12}(4), 568--581 (1964). doi: 10.1287/opre.12.4.568

\bibitem{PerronFurnon2023ORTools}
Didier, F., Perron, L., Mohajeri, S., Gay, S.A., Cuvelier, T., Furnon, V.: OR-Tools' vehicle routing solver: A generic constraint-programming solver with heuristic search for routing problems. In: Congrès de la Société française de recherche opérationnelle et d'aide à la décision (ROADEF 2023) (2023)

\bibitem{Voudouris1999GLS}
Voudouris, C., Tsang, E.: Guided local search and its application to the traveling salesman problem. Eur. J. Oper. Res. \textbf{113}(2), 469--499 (1999). doi: 10.1016/S0377-2217(98)00099-X

\bibitem{Holland1992Adaptation}
Holland, J.H.: Adaptation in Natural and Artificial Systems: An Introductory Analysis with Applications to Biology, Control, and Artificial Intelligence. MIT Press, Cambridge, MA (1992). doi: 10.7551/mitpress/1090.001.0001

\bibitem{Joshi2019GCN}
Joshi, C.K., Laurent, T., Bresson, X.: An efficient graph convolutional network technique for the travelling salesman problem. arXiv:1906.01227 (2019)

\bibitem{Dahiya2018TSPReview}
Dahiya, C., Sangwan, S.: Literature review on travelling salesman problem. Int. J. Res. \textbf{5}(16), 1152--1155 (2018)

\bibitem{Wu2024NCOSurvey}
Wu, X., Wang, D., Wen, L., Xiao, Y., Wu, C., Wu, Y., Yu, C., Maskell, D.L., Zhou, Y.: Neural combinatorial optimization algorithms for solving vehicle routing problems: A comprehensive survey with perspectives. arXiv:2406.00415 (2024). doi: 10.48550/arXiv.2406.00415

\bibitem{Boeing2017OSMnx}
Boeing, G.: OSMnx: New methods for acquiring, constructing, analyzing, and visualizing complex street networks. Comput. Environ. Urban Syst. \textbf{65}, 126--139 (2017). doi: 10.1016/j.compenvurbsys.2017.05.004

\bibitem{Pedregosa2011Sklearn}
Pedregosa, F., Varoquaux, G., Gramfort, A., et al.: Scikit-learn: Machine learning in Python. J. Mach. Learn. Res. \textbf{12}, 2825--2830 (2011)

\bibitem{Rousseeuw1987Silhouettes}
Rousseeuw, P.J.: Silhouettes: A graphical aid to the interpretation and validation of cluster analysis. J. Comput. Appl. Math. \textbf{20}, 53--65 (1987). doi: 10.1016/0377-0427(87)90125-7

\bibitem{Demsar2006Statistical}
Demšar, J.: Statistical comparisons of classifiers over multiple data sets. J. Mach. Learn. Res. \textbf{7}, 1--30 (2006)

\end{thebibliography}

\end{document}
